\documentclass[11pt,a4paper]{article}

\usepackage[top=2.4cm,bottom=2.4cm,left=2.6cm,right=2.6cm]{geometry}
\usepackage{amsmath,amssymb,amsthm}
\usepackage{mathtools}
\usepackage{booktabs}
\usepackage[expansion=false]{microtype}
\usepackage{enumitem}
\usepackage[font=small]{caption}
\usepackage{xcolor}
\usepackage[colorlinks=true,linkcolor=blue!45!black,citecolor=blue!45!black,urlcolor=blue]{hyperref}
\usepackage[giveninits=true,maxbibnames=10,url=false,isbn=false]{biblatex}
\addbibresource{references.bib}

\newcommand{\dd}{\mathrm{d}}
\newcommand{\Var}{\operatorname{Var}}
\newcommand{\Cov}{\operatorname{Cov}}
\newcommand{\Corr}{\operatorname{Corr}}
\newcommand{\Ai}{\operatorname{Ai}}
\newcommand{\Bi}{\operatorname{Bi}}
\newcommand{\E}{\mathbb{E}}
\DeclareMathOperator*{\argmax}{\arg\,max}
\DeclareMathOperator*{\argmin}{\arg\,min}

\title{Forecasting tipping times from data}
\author{}
\date{}

\begin{document}
\maketitle
\vspace{-1.1cm}

\begin{abstract}\noindent
        Combining recent results in nonautonomous dynamics and statistical estimation for stochastic differential equations we derive a fully data-driven pipeline for the accurate forecast of tipping times in complex systems.
\end{abstract}

\section{Introduction}\label{sec:intro}
We are concerned with the problem of forecasting the time at which complex systems will undergo critical transitions given past and present data.
The forecast is a (linear) extrapolation of the estimated return rate of nonlinear systems slowly ramped towards a saddle-node bifurcation.
In \S\ref{sec:tipping} we introduce the mathematical formalism of tipping times in slow-fast stochastic processes.
In \S\ref{sec:formulation} we formulate the extrapolation problem and present three techniques to estimate the return rate of a given sample, two of which have already been proposed in the literature and a new one which we introduce.
In \S\ref{sec:observation} we derive a series of conditions that allow for the accurate forecast of the tipping times to be considered when applying this methodology to collected measurements or experiments data.
In \S\ref{sec:application} we apply the forecasting methodology and the observation protocols of \S\ref{sec:formulation} and \S\ref{sec:observation} to simulated and real-world timeseries data.
In \S\ref{sec:conclusion} we conclude.

\printbibliography

\end{document}
